\section{Multi-Channel Approaches and Link Selection Techniques}
\label{sec:sota_multichannel}

The complementarity among technologies, including dedicated radio, cellular, and VLC, has led to growing interest in multi-channel architectures, where overall reliability is pursued through \emph{diversity}, redundancy, or dynamic link selection based on channel conditions and controller requirements \cite{hardes2019communication,ishihara2015improving,segata2016platooning}. Conceptually, the literature moves along three directions: i) parallel use of multiple interfaces with message duplication, ii) assignment of different roles to the technologies, for instance primary channel and backup, and iii) dynamic adaptation of the \emph{interface + controller} pair according to the observed communication quality. This last direction is the most interesting for safety-critical applications, but also the most complex to design.

The critical point is that \emph{switching} between interfaces cannot be treated as a purely networking problem. A change in technology modifies the effective update frequency, loss profile, and latency, and therefore the quality of the information available to the controller. If communication quality degrades, proper controller operation may require wider safety margins, for example by increasing the inter-vehicle gap, thus reducing the expected benefits of platooning; in the most critical cases, deteriorating connectivity can even compromise system stability. The literature on network/control co-design and degradation mechanisms shows precisely the need to explicitly link link quality to control modes and spacing parameters \cite{giordano2019joint,ploeg2015graceful}.

In this context, the work of Segata \emph{et al.} \cite{segata2023multi-technology} is a particularly relevant reference because it jointly addresses interface quality monitoring, \emph{fallback}/\emph{recovery} logic, and safe transitions between controllers. The system assumes the simultaneous presence of three interfaces, IEEE~802.11p, VLC, and LTE C-V2X, and the redundant transmission of control messages over all active interfaces. Link quality is monitored through PDR, estimated with a per-interface \emph{ring buffer} that exploits the comparison of sequence numbers observed over the different technologies to infer losses and failures.

On this basis, the paper introduces the \emph{Safe Autonomous Switchover Algorithm} (SafeSwitch), implemented as a distributed finite-state machine. States of type $F_i$ represent stable driving conditions with $i$ available interfaces, each associated with a given controller and a given gap; states $G_i$, instead, represent intermediate \emph{gap control} phases, in which the system progressively increases or decreases the inter-vehicle distance before completing the controller change \cite{segata2023multi-technology}. This detail is essential: the transition does not occur instantaneously, but is made dynamically compatible with platoon stability.

The experimental configuration proposed by the authors uses, as an example, a PATH controller when all three interfaces are available, which is the highest-performance state, a Ploeg-type controller in states with lower connectivity, and an ACC as the final fallback when all interfaces become unavailable \cite{segata2023multi-technology}. The mechanism is also bidirectional: it handles not only degradation, or \emph{failure}, but also recovery, with PDR thresholds and confirmation times to avoid oscillations caused by transient channel recovery. If the last interface is lost, the system still falls back to a non-cooperative mode, ACC or manual in the general formulation, thus preserving a safety strategy.

An important result of the work is that the comparison with a naive fallback, namely a direct transition from CACC to ACC, highlights marked qualitative differences: the abrupt fallback can become unsafe in the presence of severe disturbances, such as sudden braking, whereas the strategy with \emph{gap control} reduces the risk and maintains a more regular behavior \cite{segata2023multi-technology}. The authors also show a less obvious but very useful aspect for this thesis: the increase in distance imposed by the fallback can in turn degrade some technologies, especially VLC, creating a feedback coupling between safety logic and channel quality. This kind of insight emerges only in simulations that jointly model networking, control, and vehicle dynamics.

For the present thesis work, this framework is directly relevant. On the one hand, \emph{SafeSwitch} provides the reference logic to be reproduced and validated; on the other hand, the paper's multi-technology architecture was designed in a modular way and already oriented toward extension to later technologies, including 5G \cite{segata2023multi-technology}. Replacing the LTE interface with 5G through Simu5G, together with updating the toolchain made of PLEXE, Veins, and SUMO, therefore fits as a natural extension of the state of the art.
